%\chapter*{Lista de símbolos}\label{LDS}


\begin{longtable}{lcp{.7\textwidth} }
\multicolumn{3}{l}{\textbf{Símbolos orbitales}}\\
\\
$a$ & --- & Semieje mayor de la órbita.\\
$e$ & --- & Excentricidad orbital.\\
$i$ & --- & Inclinación del plano orbital respecto al plano de referencia.\\
$\Omega$ & --- & Ascensión recta del nodo ascendente (RAAN).\\
$\omega$ & --- & Argumento del perigeo.\\
$\nu$ & --- & Anomalía verdadera.\\
$r$ & --- & Distancia radial del satélite al centro del cuerpo.\\
$h$ & --- & Altitud sobre la superficie del cuerpo ($r=R+h$; $h_{\min}$: límite inferior de cada capa atmosférica).\\
$v$ & --- & Velocidad orbital.\\
$T$ & --- & Periodo orbital.\\
$t$ & --- & Tiempo.\\
$n$ & --- & Movimiento medio, $n=\sqrt{\mu/a^3}$.\\
$\varepsilon$ & --- & Energía orbital específica, $\varepsilon=-\mu/(2a)$.\\
$p$ & --- & Semi-latus rectum (parámetro de la cónica), $p=a(1-e^2)$.\\
$\mu$ & --- & Parámetro gravitatorio del cuerpo central, $\mu=GM$.\\
$G$ & --- & Constante de gravitación universal.\\
$M$ & --- & Masa del cuerpo central.\\
\\
\multicolumn{3}{l}{\textbf{Perturbaciones y dinámica}}\\
\\
$J_2$ & --- & Coeficiente de achatamiento planetario (perturbación gravitatoria dominante).\\
$R$ & --- & Radio ecuatorial del cuerpo central.\\
$g$ & --- & Aceleración de la gravedad en la superficie del cuerpo.\\
$\rho$ & --- & Densidad atmosférica ($\rho_{\text{base}}$: densidad de referencia en la base de cada capa).\\
$P$ & --- & Presión atmosférica.\\
$T$ & --- & Temperatura atmosférica (mismo símbolo que el periodo; el contexto los distingue).\\
$\bar{M}$ & --- & Masa molar media de la atmósfera.\\
$R_g$ & --- & Constante de los gases.\\
$H$ & --- & Altura de escala atmosférica, $H = R_g\,T/(\bar{M}\,g)$ ($H_i$: valor en la capa $i$).\\
$N$ & --- & Número de capas del modelo atmosférico.\\
$C_d$ & --- & Coeficiente de arrastre (adimensional).\\
$C_d A/m$ & --- & Coeficiente balístico inverso (área frontal por coeficiente de arrastre dividida por masa).\\
$A$ & --- & Área frontal de la nave.\\
$m$ & --- & Masa de la nave.\\
$F_{\text{drag}}$ & --- & Fuerza de arrastre atmosférico, $F_{\text{drag}}=\tfrac{1}{2}\rho v^2 C_d A$.\\
\\
\multicolumn{3}{l}{\textbf{Maniobras y transferencias}}\\
\\
$\Delta v$ & --- & Incremento de velocidad de un impulso; coste en combustible de una maniobra ($\Delta v_1$, $\Delta v_2$: primer y segundo impulso de una transferencia).\\
$\Delta i$ & --- & Cambio de inclinación del plano orbital.\\
$C_3$ & --- & Energía característica, $C_3=v_\infty^2$ (km$^2$/s$^2$).\\
$v_\infty$ & --- & Exceso hiperbólico de velocidad (velocidad relativa al cuerpo lejos de su influencia).\\
$q$ & --- & Presión dinámica, $q=\tfrac{1}{2}\rho v^2$.\\
$q_{\max}$ & --- & Presión dinámica máxima admisible por la nave (límite estructural/térmico).\\
$R$ & --- & Razón de radios $r_2/r_1$ de una transferencia (adimensional; distíngase del radio del cuerpo).\\
$\alpha$ & --- & Ángulo girado por un impulso en una maniobra con cambio de plano.\\
$r_1,\,r_2$ & --- & Radios de las órbitas circular inicial y final de una transferencia.\\
$a_t$ & --- & Semieje de la elipse de transferencia, $a_t=(r_1+r_2)/2$.\\
$v_{c1}$ & --- & Velocidad circular de la órbita inicial, $v_{c1}=\sqrt{\mu/r_1}$ (referencia de la adimensionalización).\\
\end{longtable}
